The reconstructed spectra of Fig.~\ref{fig:spectra} are best read not as a
rate prediction but as a map of \emph{where}, on the axis a
bandwidth-limited QPD readout actually records, the reactor signal and each
ambient background come to rest. That map is the deliverable of this work:
it tells a would-be experiment which instrumental knob is decisive before
any exposure is accumulated. Its interpretation, however, is bounded by four
modeling uncertainties of very different character and provenance, and an
honest reading requires keeping them separate.

\emph{Saturated-regime response shape.} The most consequential uncertainty
is also the one with the least external support. Between the linear regime
($E_{\rm rec}\approx0.5\,E_{\rm dep}$) and the fully saturated plateau, the
shape of $R(E_{\rm rec}\,|\,E_{\rm dep})$ has no published anchor at any
energy; it is fixed only by its two limiting cases and by calibration
consistency (Sec.~\ref{sec:response}). The direct consequence is that the
tens-of-keV muon pile-up of Fig.~\ref{fig:spectra} is an \emph{instrumental
artifact} of the modeled 25~kHz ceiling rather than a physical spectral
line. What would tighten it is unambiguous: a direct QPD saturation
measurement on a thin wafer, or a full microphysical readout simulation
resolving the tunneling-burst statistics through the censoring transition.
Crucially, this caveat is quarantined to the channels that actually
saturate---muons and the high-recoil gamma tail---because the reactor-CEvNS
recoils sit entirely within the linear regime and never test the
unbenchmarked shape.

\emph{Environmental-gamma normalization.} The absolute Compton rate is a
representative, site-dependent estimate good to a factor of $\sim$2, driven
by the assumed radiogenic flux and, in particular, by the ${}^{238}$U-chain
balance rather than a measured line intensity (Sec.~\ref{sec:backgrounds}).
The Compton edge positions, by contrast, are fixed by kinematics and carry
high confidence. A factor-of-two shift rescales the continuum but preserves
both the edge structure and the qualitative conclusion that the gamma
continuum overlaps the signal region, so the \emph{ranking} of this
background is robust even though its normalization is not. Only an on-site
gamma survey at the intended deployment can replace the representative
estimate with a measured one.

\emph{Sub-1.8~MeV reactor flux.} Below the reach of the Huber--Mueller
conversion fits, the antineutrino spectrum is a cited but not directly
reproduced model placeholder, carried with a wide band that translates into
a $\sim$6--10\% uncertainty on the reconstructed CEvNS rate below $\sim$95~eV
(Sec.~\ref{sec:cevns}). This is the leading \emph{signal-side} systematic,
and it acts precisely where the reconstructed CEvNS spectrum peaks. It is,
however, a band on the rate, not an ambiguity in the spectral placement: the
signal remains anchored at tens of eV regardless of where in the band the
true low-energy flux lies. Digitizing a per-isotope summation or CONFLUX
table below 1.8~MeV, following the approach of Ref.~\cite{kopeikin2012},
would collapse this band to the percent level.

\emph{Crossover deposit energy.} The linear-to-saturated crossover rides on
the two low-confidence sensor-sharing parameters $f_{\rm prompt}$ and $r$,
for which no thin-wafer QPD measurement exists; we therefore report it as a
band spanning $\sim$50~eV to $\sim$13~keV rather than a single number
(Sec.~\ref{sec:response}). Because the CEvNS recoils fall below even the
lowest end of this band, the signal reconstructs linearly across the entire
scan, so its placement is insensitive to where the true crossover sits; the
uncertainty is inherited only by the saturated backgrounds, whose pile-up
energies shift with the band. A dedicated thin-wafer sensor-sharing
measurement would fix $f_{\rm prompt}$ and $r$ and turn the band into a line.

Read together, three of the four uncertainties act on the saturated
backgrounds and leave the flagship result---that reactor CEvNS reconstructs
to tens of eV---as the most robust feature of the analysis. This is the
sense in which the reconstructed axis is informative despite the model risk.

The implications for a QPD reactor-CEvNS search follow directly and are
mostly about detector design rather than statistics. First, \emph{threshold
placement is decisive}: the CEvNS signal populates the very lowest
reconstructed energies, so recovering it demands an eV-scale reconstruction
floor, not merely a low deposited threshold---a requirement that places the
signal below the sub-100~eV regimes now being probed by cryogenic
low-threshold detectors~\cite{nucleus} and reactor-CEvNS experiments
generally. Second, cosmic muons do \emph{not} preclude quiescent operation:
their pile-up lands at tens of keV, well \emph{above} the CEvNS band, and the
pileup occupancy is tiny ($R\tau\approx3\times10^{-5}$), so muon events are
sparse and spectrally displaced from the signal. A surface deployment would
nonetheless benefit from muon tagging or an active veto, since the rare
long-chord muon tail is the very driver of readout saturation and any
misreconstruction leaks downward in energy. Third, and most limiting for a
signal search, the environmental-gamma Compton continuum \emph{overlaps} the
CEvNS region, at the level of $\sim$$10^{3}$ counts~kg$^{-1}$day$^{-1}$
keV$^{-1}$; unlike the muons it sits directly on top of the signal and is the
dominant reducible background. Its mitigation---passive shielding and,
ultimately, underground or low-background siting---is therefore the primary
lever for turning this forward model into a viable search. These conclusions
place the QPD concept squarely within the reactor-CEvNS program pioneered at
the ionization scale by CONUS+~\cite{conus2025}, the Dresden-II point-contact
germanium result~\cite{dresden2022}, and the germanium rate benchmarks
of Ref.~\cite{billard2017}; the distinction here is the unified phonon scale
and the bandwidth-limited reconstruction, both inherited from the QPD device
concept of Ref.~\cite{ramanathan2026}.

We close by stating the scope plainly. This is a stage-one forward-model
feasibility study: it establishes the reconstructed-energy regime in which a
QPD germanium wafer would operate and identifies the design choices that
regime forces. It does \emph{not} compute a sensitivity, and it makes no
claim of detectability or discovery significance. Converting this map into a
sensitivity would require an exposure model, a noise-and-resolution
treatment of the reconstruction floor, and---above all---the measured inputs
that the four caveats above call for: a QPD saturation calibration, an
on-site gamma survey, a digitized low-energy reactor flux, and a thin-wafer
sensor-sharing measurement. Each is a concrete, bounded next step, and
together they define the path from the present feasibility map to a
quantitative reactor-CEvNS reach.
